\documentclass[11pt,a4paper]{article}
\usepackage[utf8]{inputenc}
\usepackage[T1]{fontenc}
\usepackage[spanish,es-minimal,es-noshorthands]{babel}
\usepackage{lmodern}
\usepackage[sfdefault,lining]{FiraSans}
\usepackage[a4paper,top=2.3cm,bottom=2.2cm,left=1.8cm,right=1.8cm,headheight=1.3cm,headsep=0.4cm,footskip=1cm]{geometry}
\usepackage{amsmath,amssymb}
\usepackage{icomma}
\usepackage{graphicx}
\usepackage[table]{xcolor}
\usepackage{tikz}
\usepackage{fontawesome5}
\usepackage{ccicons}
\usepackage{fancyhdr}
\usepackage{lastpage}
\usepackage{tasks}
\usepackage{needspace}
\usepackage[most]{tcolorbox}
\usepackage{array}
\usepackage[hidelinks]{hyperref}
\AtBeginDocument{\hypersetup{pdftitle={Ficha de ejercicios de la unidad 6: Sistemas de ecuaciones. 3.º de ESO},pdfsubject={Matemáticas 3.º ESO},pdfkeywords={sistemas de ecuaciones, ecuaciones lineales con dos incógnitas, sustitución, igualación, reducción, método gráfico, clasificación de sistemas, problemas}}}

\definecolor{navy}{HTML}{1B2748}
\definecolor{accent}{HTML}{E8603A}
\definecolor{teal}{HTML}{1F9E90}
\definecolor{soft}{HTML}{EEF1F8}
\definecolor{muted}{HTML}{5D6685}

\setlength{\parindent}{0pt}
\setlength{\parskip}{3pt}
\settasks{label=\alph*), label-format=\bfseries\color{navy}, label-width=1.3em, item-indent=1.8em, column-sep=0.6em, after-item-skip=5pt, before-skip=3pt, after-skip=4pt}

% ---------- capçalera i peu ----------
\newcommand{\llicencia}{\href{https://creativecommons.org/licenses/by-nc-sa/4.0/deed.es}{\ccbyncsa}\ \ Licencia Creative Commons Atribución-NoComercial-CompartirIgual 4.0 Internacional (CC BY-NC-SA 4.0)}
\pagestyle{fancy}
\fancyhf{}
\fancyhead[L]{\raisebox{-0.35\height}{\includegraphics[height=0.95cm]{logo-gva.png}}}
\fancyhead[C]{\small\color{muted}Ficha de ejercicios · Unidad 6: Sistemas de ecuaciones · 3.º de ESO}
\fancyhead[R]{\raisebox{-0.35\height}{\includegraphics[height=0.95cm]{logo-vermellar.png}}}
\fancyfoot[L]{\scriptsize\color{muted}\llicencia}
\fancyfoot[R]{\small\color{muted}\thepage}
\renewcommand{\headrulewidth}{0.4pt}
\renewcommand{\headrule}{\color{muted!40}\hrule height \headrulewidth}
\fancypagestyle{primera}{\fancyhf{}\renewcommand{\headrulewidth}{0pt}
  \fancyfoot[L]{\scriptsize\color{muted}\llicencia}
  \fancyfoot[R]{\small\color{muted}\thepage}}

% ---------- nivells i exercicis ----------
\newcommand{\estrelles}[1]{{\color{accent}\ifcase#1\or\faStar\faStar[regular]\faStar[regular]\or\faStar\faStar\faStar[regular]\or\faStar\faStar\faStar\fi}}
\newcounter{exer}
\newcommand{\ex}[1]{\par\needspace{4\baselineskip}\medskip\stepcounter{exer}\noindent{\color{navy}\bfseries\theexer.}\ \ {\small\estrelles{#1}}\ \ }

\newcommand{\seccio}[2]{\par\needspace{10\baselineskip}\bigskip
  \begin{tcolorbox}[colback=navy,colframe=navy,coltext=white,arc=2mm,boxrule=0pt,left=3mm,right=3mm,top=1.2mm,bottom=1.2mm]
  {\large\bfseries #1\quad #2}
  \end{tcolorbox}\vspace{-2pt}}

\newcolumntype{C}[1]{>{\centering\arraybackslash}p{#1}}
\usepackage{pifont}

% quadre de recordatori (només a la fitxa d'exercicis)
\newcommand{\recorda}[1]{\par\smallskip\begin{tcolorbox}[colback=accent!7,colframe=accent!55,arc=2mm,boxrule=0.5pt,left=2.5mm,right=2.5mm,top=1mm,bottom=1mm]\small{\color{accent}\faLightbulb[regular]}\ \textbf{Recuerda:}\ #1\end{tcolorbox}}
\newcommand{\res}[1]{#1}
\newcommand{\rest}[1]{#1}

\def\fmneed{14}
% =====================================================================
% Capçalera FentMates de les fitxes: un color per a cada curs.
% Els colors són els de la pàgina inicial (LEVELS a assets/js/materials.js);
% el logo ve de _eines/logo/pdf/ (python3 _eines/logo/genera.py).
%
% Ús, al final del preàmbul de la unitat:
%   % =====================================================================
% Capçalera FentMates de les fitxes: un color per a cada curs.
% Els colors són els de la pàgina inicial (LEVELS a assets/js/materials.js);
% el logo ve de _eines/logo/pdf/ (python3 _eines/logo/genera.py).
%
% Ús, al final del preàmbul de la unitat:
%   % =====================================================================
% Capçalera FentMates de les fitxes: un color per a cada curs.
% Els colors són els de la pàgina inicial (LEVELS a assets/js/materials.js);
% el logo ve de _eines/logo/pdf/ (python3 _eines/logo/genera.py).
%
% Ús, al final del preàmbul de la unitat:
%   \input{../comu/fentmates.tex}
%   \fmcurs{1eso}
% i, al començament del document:
%   \capcalera{Matemàtiques}{1r ESO}{Fitxa d'exercicis · Unitat 2}{Divisibilitat}
% Opcional abans de \input: \def\fmneed{14} (espai mínim abans de cada secció).
% =====================================================================
\makeatletter
\def\fm@logodir{../../logo/pdf}
\providecommand{\fmneed}{10}
\def\fm@cursnom{}

% clau del curs, color i color del text damunt del color
\newcommand{\fm@curs}[3]{\@namedef{fm@color@#1}{#2}\@namedef{fm@text@#1}{#3}}
\fm@curs{1eso}{E6B84A}{navy}
\fm@curs{2eso}{3B82F6}{white}
\fm@curs{3eso}{FF7A50}{white}
\fm@curs{4eso-a}{A78BFA}{white}
\fm@curs{4eso-b}{6D28D9}{white}
\fm@curs{1batx-mi}{9B1C1C}{white}
\fm@curs{1batx-ccss}{EF6B6B}{white}
\fm@curs{2batx-mii}{1E6B47}{white}
\fm@curs{2batx-ccss}{2BB3A3}{white}

\newcommand{\fmcurs}[1]{%
  \@ifundefined{fm@color@#1}{\PackageError{fentmates}{Curs desconegut: #1}{}}{}%
  \edef\fm@tmp{\noexpand\definecolor{curs}{HTML}{\@nameuse{fm@color@#1}}%
               \noexpand\colorlet{curstext}{\@nameuse{fm@text@#1}}}\fm@tmp
  \colorlet{cursfosc}{curs!72!black}%
  \def\fm@logo{\fm@logodir/vertical-#1.pdf}%
  \def\fm@marca{\fm@logodir/marca-#1.pdf}}

% ---------- capçalera de la primera pàgina ----------
% #1 matèria (petit), #2 curs, #3 tipus de document i unitat, #4 títol
\newcommand{\capcalera}[4]{%
  \gdef\fm@cursnom{#2}%
  \par\noindent
  \begin{tcolorbox}[colback=curs,colframe=curs,boxrule=0pt,arc=2.2mm,width=3.5cm,height=2.15cm,
      valign=center,left=3mm,right=2mm,top=1mm,bottom=1mm,nobeforeafter,box align=center]
    \color{curstext}{\footnotesize #1}\par\vspace{1pt}{\fontsize{19}{21}\selectfont\bfseries #2\par}%
  \end{tcolorbox}\hspace{5mm}%
  \begin{minipage}[c]{\dimexpr\linewidth-3.5cm-5mm-2.9cm\relax}
    {\color{cursfosc}\small\bfseries\MakeUppercase{#3}}\par\vspace{2pt}
    {\color{navy}\fontsize{20}{23}\selectfont\bfseries #4\par}
  \end{minipage}\hfill
  \begin{minipage}[c]{2.6cm}\raggedleft\includegraphics[width=2.6cm]{\fm@logo}\end{minipage}\par
  \vspace{2pt}}

% ---------- barres de secció del color del curs ----------
\renewcommand{\seccio}[2]{\par\needspace{\fmneed\baselineskip}\bigskip
  \begin{tcolorbox}[colback=curs,colframe=curs,coltext=curstext,arc=2mm,boxrule=0pt,left=3mm,right=3mm,top=1.2mm,bottom=1.2mm]
  {\large\bfseries #1\quad #2}
  \end{tcolorbox}\vspace{-2pt}}

% ---------- estrelles de dificultat del color del curs (a Batxillerat no se'n posen) ----------
\@ifundefined{estrelles}{}{%
  \renewcommand{\estrelles}[1]{{\color{curs}\ifcase#1\or\faStar\faStar[regular]\faStar[regular]\or\faStar\faStar\faStar[regular]\or\faStar\faStar\faStar\fi}}}

% ---------- capçalera i peu de les pàgines ----------
% primera pàgina: logos institucionals al peu
\fancypagestyle{primera}{\fancyhf{}\renewcommand{\headrulewidth}{0pt}
  \fancyfoot[L]{\raisebox{-0.32\height}{\includegraphics[height=0.78cm]{logo-gva.png}}\hspace{5mm}%
    \raisebox{-0.32\height}{\includegraphics[height=0.78cm]{logo-vermellar.png}}}
  \fancyfoot[R]{\scriptsize\color{muted}\href{https://creativecommons.org/licenses/by-nc-sa/4.0/deed.ca}{\ccbyncsa}\ \ CC BY-NC-SA 4.0 · fentmates.org\qquad\small\thepage}}
% la resta: el símbol a l'esquerra i el curs a la dreta, amb el títol al centre com abans
\fancyhead[L]{\raisebox{-0.3\height}{\includegraphics[height=0.62cm]{\fm@marca}}}
\fancyhead[R]{\small\bfseries\color{cursfosc}\fm@cursnom}
\renewcommand{\headrule}{\color{curs!45}\hrule height \headrulewidth}
\makeatother

%   \fmcurs{1eso}
% i, al començament del document:
%   \capcalera{Matemàtiques}{1r ESO}{Fitxa d'exercicis · Unitat 2}{Divisibilitat}
% Opcional abans de \input: \def\fmneed{14} (espai mínim abans de cada secció).
% =====================================================================
\makeatletter
\def\fm@logodir{../../logo/pdf}
\providecommand{\fmneed}{10}
\def\fm@cursnom{}

% clau del curs, color i color del text damunt del color
\newcommand{\fm@curs}[3]{\@namedef{fm@color@#1}{#2}\@namedef{fm@text@#1}{#3}}
\fm@curs{1eso}{E6B84A}{navy}
\fm@curs{2eso}{3B82F6}{white}
\fm@curs{3eso}{FF7A50}{white}
\fm@curs{4eso-a}{A78BFA}{white}
\fm@curs{4eso-b}{6D28D9}{white}
\fm@curs{1batx-mi}{9B1C1C}{white}
\fm@curs{1batx-ccss}{EF6B6B}{white}
\fm@curs{2batx-mii}{1E6B47}{white}
\fm@curs{2batx-ccss}{2BB3A3}{white}

\newcommand{\fmcurs}[1]{%
  \@ifundefined{fm@color@#1}{\PackageError{fentmates}{Curs desconegut: #1}{}}{}%
  \edef\fm@tmp{\noexpand\definecolor{curs}{HTML}{\@nameuse{fm@color@#1}}%
               \noexpand\colorlet{curstext}{\@nameuse{fm@text@#1}}}\fm@tmp
  \colorlet{cursfosc}{curs!72!black}%
  \def\fm@logo{\fm@logodir/vertical-#1.pdf}%
  \def\fm@marca{\fm@logodir/marca-#1.pdf}}

% ---------- capçalera de la primera pàgina ----------
% #1 matèria (petit), #2 curs, #3 tipus de document i unitat, #4 títol
\newcommand{\capcalera}[4]{%
  \gdef\fm@cursnom{#2}%
  \par\noindent
  \begin{tcolorbox}[colback=curs,colframe=curs,boxrule=0pt,arc=2.2mm,width=3.5cm,height=2.15cm,
      valign=center,left=3mm,right=2mm,top=1mm,bottom=1mm,nobeforeafter,box align=center]
    \color{curstext}{\footnotesize #1}\par\vspace{1pt}{\fontsize{19}{21}\selectfont\bfseries #2\par}%
  \end{tcolorbox}\hspace{5mm}%
  \begin{minipage}[c]{\dimexpr\linewidth-3.5cm-5mm-2.9cm\relax}
    {\color{cursfosc}\small\bfseries\MakeUppercase{#3}}\par\vspace{2pt}
    {\color{navy}\fontsize{20}{23}\selectfont\bfseries #4\par}
  \end{minipage}\hfill
  \begin{minipage}[c]{2.6cm}\raggedleft\includegraphics[width=2.6cm]{\fm@logo}\end{minipage}\par
  \vspace{2pt}}

% ---------- barres de secció del color del curs ----------
\renewcommand{\seccio}[2]{\par\needspace{\fmneed\baselineskip}\bigskip
  \begin{tcolorbox}[colback=curs,colframe=curs,coltext=curstext,arc=2mm,boxrule=0pt,left=3mm,right=3mm,top=1.2mm,bottom=1.2mm]
  {\large\bfseries #1\quad #2}
  \end{tcolorbox}\vspace{-2pt}}

% ---------- estrelles de dificultat del color del curs (a Batxillerat no se'n posen) ----------
\@ifundefined{estrelles}{}{%
  \renewcommand{\estrelles}[1]{{\color{curs}\ifcase#1\or\faStar\faStar[regular]\faStar[regular]\or\faStar\faStar\faStar[regular]\or\faStar\faStar\faStar\fi}}}

% ---------- capçalera i peu de les pàgines ----------
% primera pàgina: logos institucionals al peu
\fancypagestyle{primera}{\fancyhf{}\renewcommand{\headrulewidth}{0pt}
  \fancyfoot[L]{\raisebox{-0.32\height}{\includegraphics[height=0.78cm]{logo-gva.png}}\hspace{5mm}%
    \raisebox{-0.32\height}{\includegraphics[height=0.78cm]{logo-vermellar.png}}}
  \fancyfoot[R]{\scriptsize\color{muted}\href{https://creativecommons.org/licenses/by-nc-sa/4.0/deed.ca}{\ccbyncsa}\ \ CC BY-NC-SA 4.0 · fentmates.org\qquad\small\thepage}}
% la resta: el símbol a l'esquerra i el curs a la dreta, amb el títol al centre com abans
\fancyhead[L]{\raisebox{-0.3\height}{\includegraphics[height=0.62cm]{\fm@marca}}}
\fancyhead[R]{\small\bfseries\color{cursfosc}\fm@cursnom}
\renewcommand{\headrule}{\color{curs!45}\hrule height \headrulewidth}
\makeatother

%   \fmcurs{1eso}
% i, al començament del document:
%   \capcalera{Matemàtiques}{1r ESO}{Fitxa d'exercicis · Unitat 2}{Divisibilitat}
% Opcional abans de \input: \def\fmneed{14} (espai mínim abans de cada secció).
% =====================================================================
\makeatletter
\def\fm@logodir{../../logo/pdf}
\providecommand{\fmneed}{10}
\def\fm@cursnom{}

% clau del curs, color i color del text damunt del color
\newcommand{\fm@curs}[3]{\@namedef{fm@color@#1}{#2}\@namedef{fm@text@#1}{#3}}
\fm@curs{1eso}{E6B84A}{navy}
\fm@curs{2eso}{3B82F6}{white}
\fm@curs{3eso}{FF7A50}{white}
\fm@curs{4eso-a}{A78BFA}{white}
\fm@curs{4eso-b}{6D28D9}{white}
\fm@curs{1batx-mi}{9B1C1C}{white}
\fm@curs{1batx-ccss}{EF6B6B}{white}
\fm@curs{2batx-mii}{1E6B47}{white}
\fm@curs{2batx-ccss}{2BB3A3}{white}

\newcommand{\fmcurs}[1]{%
  \@ifundefined{fm@color@#1}{\PackageError{fentmates}{Curs desconegut: #1}{}}{}%
  \edef\fm@tmp{\noexpand\definecolor{curs}{HTML}{\@nameuse{fm@color@#1}}%
               \noexpand\colorlet{curstext}{\@nameuse{fm@text@#1}}}\fm@tmp
  \colorlet{cursfosc}{curs!72!black}%
  \def\fm@logo{\fm@logodir/vertical-#1.pdf}%
  \def\fm@marca{\fm@logodir/marca-#1.pdf}}

% ---------- capçalera de la primera pàgina ----------
% #1 matèria (petit), #2 curs, #3 tipus de document i unitat, #4 títol
\newcommand{\capcalera}[4]{%
  \gdef\fm@cursnom{#2}%
  \par\noindent
  \begin{tcolorbox}[colback=curs,colframe=curs,boxrule=0pt,arc=2.2mm,width=3.5cm,height=2.15cm,
      valign=center,left=3mm,right=2mm,top=1mm,bottom=1mm,nobeforeafter,box align=center]
    \color{curstext}{\footnotesize #1}\par\vspace{1pt}{\fontsize{19}{21}\selectfont\bfseries #2\par}%
  \end{tcolorbox}\hspace{5mm}%
  \begin{minipage}[c]{\dimexpr\linewidth-3.5cm-5mm-2.9cm\relax}
    {\color{cursfosc}\small\bfseries\MakeUppercase{#3}}\par\vspace{2pt}
    {\color{navy}\fontsize{20}{23}\selectfont\bfseries #4\par}
  \end{minipage}\hfill
  \begin{minipage}[c]{2.6cm}\raggedleft\includegraphics[width=2.6cm]{\fm@logo}\end{minipage}\par
  \vspace{2pt}}

% ---------- barres de secció del color del curs ----------
\renewcommand{\seccio}[2]{\par\needspace{\fmneed\baselineskip}\bigskip
  \begin{tcolorbox}[colback=curs,colframe=curs,coltext=curstext,arc=2mm,boxrule=0pt,left=3mm,right=3mm,top=1.2mm,bottom=1.2mm]
  {\large\bfseries #1\quad #2}
  \end{tcolorbox}\vspace{-2pt}}

% ---------- estrelles de dificultat del color del curs (a Batxillerat no se'n posen) ----------
\@ifundefined{estrelles}{}{%
  \renewcommand{\estrelles}[1]{{\color{curs}\ifcase#1\or\faStar\faStar[regular]\faStar[regular]\or\faStar\faStar\faStar[regular]\or\faStar\faStar\faStar\fi}}}

% ---------- capçalera i peu de les pàgines ----------
% primera pàgina: logos institucionals al peu
\fancypagestyle{primera}{\fancyhf{}\renewcommand{\headrulewidth}{0pt}
  \fancyfoot[L]{\raisebox{-0.32\height}{\includegraphics[height=0.78cm]{logo-gva.png}}\hspace{5mm}%
    \raisebox{-0.32\height}{\includegraphics[height=0.78cm]{logo-vermellar.png}}}
  \fancyfoot[R]{\scriptsize\color{muted}\href{https://creativecommons.org/licenses/by-nc-sa/4.0/deed.ca}{\ccbyncsa}\ \ CC BY-NC-SA 4.0 · fentmates.org\qquad\small\thepage}}
% la resta: el símbol a l'esquerra i el curs a la dreta, amb el títol al centre com abans
\fancyhead[L]{\raisebox{-0.3\height}{\includegraphics[height=0.62cm]{\fm@marca}}}
\fancyhead[R]{\small\bfseries\color{cursfosc}\fm@cursnom}
\renewcommand{\headrule}{\color{curs!45}\hrule height \headrulewidth}
\makeatother

\fmcurs{3eso}

\begin{document}
\thispagestyle{primera}

\capcalera{Matemáticas}{3.º ESO}{Ficha de ejercicios · Unidad 6}{Sistemas de ecuaciones}
\vspace{10pt}
Nombre y apellidos: \rule{6.6cm}{0.4pt}\hfill Grupo: \rule{1.6cm}{0.4pt}\hfill Fecha: \rule{2.2cm}{0.4pt}
\vspace{8pt}

\begin{tcolorbox}[colback=soft,colframe=soft,arc=2mm,left=3mm,right=3mm,top=1.5mm,bottom=1.5mm]
\small Ejercicios graduados por niveles de dificultad:\quad \estrelles{1}\ básico \qquad \estrelles{2}\ medio \qquad \estrelles{3}\ avanzado.\\[2pt] Resuélvelos en el cuaderno mostrando todos los pasos y comprueba siempre la solución en las dos ecuaciones. Para el método gráfico, utiliza papel cuadriculado.
\end{tcolorbox}

% =====================================================
\seccio{1}{Ecuaciones lineales con dos incógnitas}

\ex{1} ¿Cuáles de estas ecuaciones son lineales con dos incógnitas?\begin{tasks}(3)
  \task $3x-y=2$
  \task $x^2+y=1$
  \task $xy=4$
  \task $\dfrac{x}{2}+y=3$
  \task $2x+5=3x$
  \task $y=4x-1$
\end{tasks}

\ex{1} Comprueba cuáles de estos pares son solución de $2x+3y=12$:\begin{tasks}(4)
  \task $(3,\,2)$
  \task $(0,\,4)$
  \task $(6,\,1)$
  \task $(-3,\,6)$
\end{tasks}

\ex{1} Despeja la $y$ y completa una tabla de valores con $x=-2,\,-1,\,0,\,1,\,2$:\begin{tasks}(2)
  \task $x+y=5$
  \task $2x-y=1$
\end{tasks}\recorda{Cada par $(x,\,y)$ de la tabla es una solución de la ecuación. Una ecuación con dos incógnitas tiene infinitas soluciones.}

\ex{2} Escribe tres soluciones de cada ecuación:\begin{tasks}(2)
  \task $x-3y=6$
  \task $4x+2y=10$
\end{tasks}

\ex{2} Responde:\begin{tasks}(1)
  \task Halla $k$ para que $(-1,\,3)$ sea solución de $4x+ky=8$.
  \task Halla $m$ para que $(2,\,m)$ sea solución de $5x-2y=4$.
\end{tasks}

\ex{2} Haz una tabla de valores y representa cada recta en unos ejes:\begin{tasks}(2)
  \task $x+2y=4$
  \task $3x-y=3$
\end{tasks}

\ex{3} Responde:\begin{tasks}(1)
  \task Halla la solución de $3x+2y=21$ en la que la $y$ vale 3 unidades más que la $x$.
  \task ¿En qué puntos corta la recta $x-2y=6$ a los ejes de coordenadas?
\end{tasks}


% =====================================================
\seccio{2}{Sistemas de ecuaciones y soluciones}

\ex{1} Comprueba si el par $(2,\,-1)$ es solución de cada sistema:\begin{tasks}(2)
  \task $\begin{cases}x+y=1\\ 2x-y=5\end{cases}$
  \task $\begin{cases}3x+y=5\\ x-2y=3\end{cases}$
\end{tasks}\recorda{Un par es solución del sistema si cumple \textbf{las dos} ecuaciones a la vez.}

\ex{1} ¿Cuál de estos pares es la solución del sistema $\begin{cases}x+y=3\\ 2x-y=3\end{cases}$? \quad $(1,\,2)$, \ $(2,\,1)$, \ $(3,\,0)$

\ex{2} Escribe un sistema de dos ecuaciones que tenga por solución $x=-2,\ y=5$.

\ex{3} Halla $a$ y $b$ para que $(3,\,-2)$ sea la solución del sistema $\begin{cases}ax+2y=5\\ 4x+by=6\end{cases}$.


% =====================================================
\seccio{3}{Método de sustitución}

\ex{1} Resuelve por sustitución (una incógnita ya está despejada):\begin{tasks}(2)
  \task $\begin{cases}x=3y\\ x+y=8\end{cases}$
  \task $\begin{cases}y=2x+1\\ 3x+y=11\end{cases}$
\end{tasks}

\ex{1} Resuelve por sustitución:\begin{tasks}(2)
  \task $\begin{cases}x+y=7\\ 2x-y=2\end{cases}$
  \task $\begin{cases}x-y=1\\ 3x+2y=13\end{cases}$
\end{tasks}\recorda{Despeja la incógnita que tiene coeficiente 1 o $-1$ y sustitúyela, \textbf{entre paréntesis}, en la otra ecuación.}

\ex{2} Resuelve por sustitución:\begin{tasks}(2)
  \task $\begin{cases}3x+y=1\\ 5x-2y=20\end{cases}$
  \task $\begin{cases}2x-y=-4\\ 4x+3y=2\end{cases}$
\end{tasks}

\ex{2} Resuelve por sustitución y comprueba la solución:\begin{tasks}(2)
  \task $\begin{cases}x+4y=-5\\ 3x-2y=13\end{cases}$
  \task $\begin{cases}5x-y=13\\ 2x+3y=12\end{cases}$
\end{tasks}

\ex{3} Resuelve por sustitución (ningún coeficiente es 1: aparecerán fracciones):\begin{tasks}(2)
  \task $\begin{cases}2x+3y=12\\ 5x-4y=7\end{cases}$
  \task $\begin{cases}3x-2y=-1\\ 4x+5y=14\end{cases}$
\end{tasks}

\ex{3} Al resolver $\begin{cases}y=3-x\\ 2x-3y=1\end{cases}$ un alumno escribe $2x-9-3x=1$. ¿Qué error ha cometido? Resuelve el sistema correctamente.


% =====================================================
\seccio{4}{Método de igualación}

\ex{1} Resuelve por igualación:\begin{tasks}(2)
  \task $\begin{cases}y=x+5\\ y=3x-1\end{cases}$
  \task $\begin{cases}x=2y-3\\ x=y+1\end{cases}$
\end{tasks}\recorda{Despeja \textbf{la misma} incógnita en las dos ecuaciones e iguala las expresiones.}

\ex{1} Resuelve por igualación:\begin{tasks}(2)
  \task $\begin{cases}x+y=6\\ x-y=2\end{cases}$
  \task $\begin{cases}y-2x=1\\ y+x=7\end{cases}$
\end{tasks}

\ex{2} Resuelve por igualación:\begin{tasks}(2)
  \task $\begin{cases}2x+y=8\\ 3x+y=11\end{cases}$
  \task $\begin{cases}x+3y=10\\ x-2y=0\end{cases}$
\end{tasks}

\ex{2} Resuelve por igualación y comprueba la solución:\begin{tasks}(2)
  \task $\begin{cases}3x-y=7\\ x+2y=0\end{cases}$
  \task $\begin{cases}4x+y=-2\\ 2x-y=-4\end{cases}$
\end{tasks}

\ex{3} Resuelve por igualación (aparecerán denominadores):\begin{tasks}(2)
  \task $\begin{cases}2x+3y=7\\ 3x-2y=4\end{cases}$
  \task $\begin{cases}5x-2y=1\\ 3x+4y=11\end{cases}$
\end{tasks}

\ex{3} Resuelve el sistema $\begin{cases}y=\dfrac{x+1}{2}\\[6pt] y=\dfrac{2x-1}{3}\end{cases}$ por el método más adecuado.


% =====================================================
\seccio{5}{Método de reducción}

\ex{1} Resuelve por reducción (los coeficientes de una incógnita ya son opuestos):\begin{tasks}(2)
  \task $\begin{cases}x+y=9\\ x-y=3\end{cases}$
  \task $\begin{cases}3x+2y=13\\ x-2y=-1\end{cases}$
\end{tasks}\recorda{Si los coeficientes de una incógnita son opuestos, al sumar las ecuaciones esa incógnita desaparece.}

\ex{1} Resuelve por reducción:\begin{tasks}(2)
  \task $\begin{cases}2x+y=7\\ 2x-3y=-5\end{cases}$
  \task $\begin{cases}5x+3y=11\\ 2x-3y=-4\end{cases}$
\end{tasks}

\ex{2} Resuelve por reducción (hay que multiplicar una de las ecuaciones):\begin{tasks}(2)
  \task $\begin{cases}x+2y=4\\ 3x-y=5\end{cases}$
  \task $\begin{cases}4x+3y=10\\ 2x+5y=12\end{cases}$
\end{tasks}

\ex{2} Resuelve por reducción (hay que multiplicar las dos ecuaciones):\begin{tasks}(2)
  \task $\begin{cases}3x+4y=18\\ 5x-2y=4\end{cases}$
  \task $\begin{cases}2x-5y=-1\\ 3x+2y=8\end{cases}$
\end{tasks}

\ex{3} Resuelve por reducción:\begin{tasks}(2)
  \task $\begin{cases}4x+3y=-1\\ 5x+2y=4\end{cases}$
  \task $\begin{cases}6x-5y=7\\ 4x-3y=5\end{cases}$
\end{tasks}

\ex{3} Resuelve el sistema $\begin{cases}2x+3y=4\\ 3x-2y=1\end{cases}$ por \textbf{doble reducción}: elimina primero la $y$ para hallar la $x$ y después elimina la $x$ para hallar la $y$. ¿Por qué es útil aquí?


% =====================================================
\seccio{6}{Método gráfico}

\ex{1} Lee la solución de cada sistema a partir del punto de corte de las dos rectas:\begin{tasks}(2)
  \task \begin{tikzpicture}[scale=0.4,baseline=(current bounding box.center)]
\draw[muted!25,very thin] (-5,-5) grid (5,5);
\draw[-latex,muted] (-5.4,0)--(5.6,0) node[right,font=\scriptsize]{$x$};
\draw[-latex,muted] (0,-5.4)--(0,5.6) node[above,font=\scriptsize]{$y$};
\node[font=\tiny,muted,below] at (-4,0) {-4}; \node[font=\tiny,muted,left] at (0,-4) {-4};
\node[font=\tiny,muted,below] at (-2,0) {-2}; \node[font=\tiny,muted,left] at (0,-2) {-2};
\node[font=\tiny,muted,below] at (2,0) {2}; \node[font=\tiny,muted,left] at (0,2) {2};
\node[font=\tiny,muted,below] at (4,0) {4}; \node[font=\tiny,muted,left] at (0,4) {4};
\begin{scope}\clip (-5,-5) rectangle (5,5);
\draw[accent,line width=1.1pt] (-6.000,10.000)--(6.000,-2.000);
\draw[teal,line width=1.1pt] (-6.000,-4.000)--(6.000,8.000);
\end{scope}
\end{tikzpicture}
  \task \begin{tikzpicture}[scale=0.4,baseline=(current bounding box.center)]
\draw[muted!25,very thin] (-5,-5) grid (5,5);
\draw[-latex,muted] (-5.4,0)--(5.6,0) node[right,font=\scriptsize]{$x$};
\draw[-latex,muted] (0,-5.4)--(0,5.6) node[above,font=\scriptsize]{$y$};
\node[font=\tiny,muted,below] at (-4,0) {-4}; \node[font=\tiny,muted,left] at (0,-4) {-4};
\node[font=\tiny,muted,below] at (-2,0) {-2}; \node[font=\tiny,muted,left] at (0,-2) {-2};
\node[font=\tiny,muted,below] at (2,0) {2}; \node[font=\tiny,muted,left] at (0,2) {2};
\node[font=\tiny,muted,below] at (4,0) {4}; \node[font=\tiny,muted,left] at (0,4) {4};
\begin{scope}\clip (-5,-5) rectangle (5,5);
\draw[accent,line width=1.1pt] (-6.000,-15.000)--(6.000,9.000);
\draw[teal,line width=1.1pt] (-6.000,2.500)--(6.000,-3.500);
\end{scope}
\end{tikzpicture}
\end{tasks}

\ex{1} Resuelve gráficamente:\begin{tasks}(2)
  \task $\begin{cases}x+y=6\\ x-y=2\end{cases}$
  \task $\begin{cases}y=x+1\\ y=-x+5\end{cases}$
\end{tasks}\recorda{Haz una tabla de valores para cada ecuación (dos o tres puntos), dibuja las dos rectas en los mismos ejes y busca el punto de corte.}

\ex{2} Resuelve gráficamente y comprueba la solución:\begin{tasks}(2)
  \task $\begin{cases}2x+y=4\\ x-y=-1\end{cases}$
  \task $\begin{cases}x+2y=0\\ 3x-y=7\end{cases}$
\end{tasks}

\ex{2} Resuelve gráficamente:\begin{tasks}(2)
  \task $\begin{cases}y=2x-4\\ x+y=-1\end{cases}$
  \task $\begin{cases}x-3y=3\\ 2x+y=-8\end{cases}$
\end{tasks}

\ex{3} Representa estos sistemas. ¿Qué observas? ¿Cuántas soluciones tienen?\begin{tasks}(2)
  \task $\begin{cases}x+y=2\\ 2x+2y=6\end{cases}$
  \task $\begin{cases}x-2y=2\\ 3x-6y=6\end{cases}$
\end{tasks}


% =====================================================
\seccio{7}{Clasificación de los sistemas}

\ex{1} Clasifica estos sistemas sin resolverlos (SCD, SCI o SI):\begin{tasks}(2)
  \task $\begin{cases}x+3y=2\\ 2x+6y=4\end{cases}$
  \task $\begin{cases}2x-y=1\\ 4x-2y=3\end{cases}$
  \task $\begin{cases}x+y=5\\ 2x-y=1\end{cases}$
  \task $\begin{cases}3x-6y=9\\ -x+2y=-3\end{cases}$
\end{tasks}\recorda{$\dfrac{a}{a'}\neq\dfrac{b}{b'}$: compatible determinado (secantes). \ $\dfrac{a}{a'}=\dfrac{b}{b'}=\dfrac{c}{c'}$: compatible indeterminado (coincidentes). \ $\dfrac{a}{a'}=\dfrac{b}{b'}\neq\dfrac{c}{c'}$: incompatible (paralelas).}

\ex{1} Completa: si las rectas de un sistema\ldots\begin{tasks}(1)
  \task se cortan en un punto, el sistema tiene \rule{2.5cm}{0.4pt} y es \rule{3.5cm}{0.4pt}.
  \task son paralelas, el sistema tiene \rule{2.5cm}{0.4pt} y es \rule{3.5cm}{0.4pt}.
  \task coinciden, el sistema tiene \rule{2.5cm}{0.4pt} y es \rule{3.5cm}{0.4pt}.
\end{tasks}

\ex{2} Resuelve por reducción y clasifica:\begin{tasks}(2)
  \task $\begin{cases}x+y=3\\ 2x+2y=5\end{cases}$
  \task $\begin{cases}x-2y=1\\ 3x-6y=3\end{cases}$
\end{tasks}

\ex{2} Resuelve y clasifica:\begin{tasks}(3)
  \task $\begin{cases}2x+3y=5\\ 4x+6y=10\end{cases}$
  \task $\begin{cases}x-y=2\\ -2x+2y=4\end{cases}$
  \task $\begin{cases}3x+y=2\\ x-y=2\end{cases}$
\end{tasks}

\ex{2} Escribe una ecuación que forme con $2x-y=3$ un sistema:\begin{tasks}(3)
  \task compatible determinado
  \task compatible indeterminado
  \task incompatible
\end{tasks}

\ex{3} ¿Para qué valor de $m$ el sistema $\begin{cases}mx+2y=3\\ 3x+y=1\end{cases}$ es incompatible? ¿Cómo es para los demás valores de $m$?

\ex{3} Halla $a$ y $b$ para que el sistema $\begin{cases}2x+ay=b\\ 4x-6y=10\end{cases}$ tenga infinitas soluciones.


% =====================================================
\seccio{8}{Sistemas con paréntesis y denominadores}

\ex{1} Escribe cada ecuación en la forma $ax+by=c$ y resuelve:\begin{tasks}(2)
  \task $\begin{cases}2(x-1)+y=7\\ x-(y-3)=3\end{cases}$
  \task $\begin{cases}3(x+y)-2x=10\\ 2(x-y)+y=-1\end{cases}$
\end{tasks}\recorda{Primero quitamos paréntesis y denominadores, transponemos y reducimos. Después elegimos el método más adecuado.}

\ex{1} Quita denominadores y resuelve:\begin{tasks}(2)
  \task $\begin{cases}\dfrac{x}{2}+y=4\\ x-\dfrac{y}{3}=1\end{cases}$
  \task $\begin{cases}\dfrac{x}{3}+\dfrac{y}{2}=2\\ \dfrac{x}{2}-\dfrac{y}{4}=1\end{cases}$
\end{tasks}

\ex{2} Resuelve:\begin{tasks}(2)
  \task $\begin{cases}\dfrac{x+1}{2}+\dfrac{y-1}{3}=3\\ x-y=-1\end{cases}$
  \task $\begin{cases}\dfrac{x-2}{3}=\dfrac{y+1}{4}\\ 2x-y=7\end{cases}$
\end{tasks}

\ex{2} Resuelve (multiplica por 10 para eliminar los decimales):\begin{tasks}(1)
  \task $\begin{cases}0{,}5x+0{,}2y=2\\ x-y=-3\end{cases}$
\end{tasks}

\ex{3} Resuelve:\begin{tasks}(1)
  \task $\begin{cases}\dfrac{2x-1}{3}-\dfrac{y+2}{2}=-1\\ 3(x-y)-2(x+1)=-6\end{cases}$
\end{tasks}

\ex{3} ¿Qué método es el más adecuado para cada sistema? Justifícalo y resuélvelos.\begin{tasks}(3)
  \task $\begin{cases}y=3x-2\\ y=x+4\end{cases}$
  \task $\begin{cases}x+5y=9\\ 2x-3y=5\end{cases}$
  \task $\begin{cases}4x+3y=6\\ 6x-3y=24\end{cases}$
\end{tasks}


% =====================================================
\seccio{9}{Problemas con sistemas de ecuaciones}

\ex{1} La suma de dos números es 56 y su diferencia es 12. ¿Cuáles son?\recorda{Elige las dos incógnitas y escribe qué representan, busca dos relaciones para plantear el sistema, resuélvelo y comprueba que la solución tiene sentido en el problema.}

\ex{1} Un número es el triple de otro y entre los dos suman 48. ¿Cuáles son?

\ex{1} En un corral hay gallinas y conejos. En total hay 40 cabezas y 112 patas. ¿Cuántos animales hay de cada tipo?

\ex{1} 3 kg de manzanas y 2 kg de peras cuestan 8 €, y 1 kg de manzanas y 3 kg de peras cuestan 5 €. ¿Cuánto cuesta el kilo de cada fruta?

\ex{2} En una hucha hay 22 monedas, unas de 2 € y otras de 50 céntimos, que suman 26 €. ¿Cuántas hay de cada tipo?

\ex{2} La entrada de un cine cuesta 7 € para los adultos y 4 € para los niños. Un día se vendieron 85 entradas y se recaudaron 445 €. ¿Cuántas entradas de cada tipo se vendieron?

\ex{2} La edad de una madre es el cuádruple de la de su hijo. Dentro de 5 años solo será el triple. ¿Cuántos años tienen ahora?

\ex{2} El perímetro de un rectángulo mide 46 m y la base mide 7 m más que la altura. ¿Qué dimensiones tiene?

\ex{2} Un test tiene 30 preguntas y las has contestado todas. Cada acierto vale 2 puntos y cada fallo resta 1. Si has sacado 39 puntos, ¿cuántas has acertado?

\ex{2} En una clase hay 28 alumnos. Si llegaran 2 chicas más, habría tantas chicas como chicos. ¿Cuántos chicos y cuántas chicas hay?

\ex{3} Queremos 60 kg de café de 7 €/kg mezclando café de 9 €/kg con café de 6 €/kg. ¿Cuántos kilos de cada tipo tenemos que poner?

\ex{3} Las cifras de un número de dos cifras suman 11. Si las intercambiamos, el número que resulta es 27 unidades menor. ¿Cuál es el número?

\ex{3} Dos ciudades están a 75 km. Dos ciclistas salen a la vez, cada uno de una ciudad, y van el uno hacia el otro. Uno va 5 km/h más rápido que el otro y se encuentran al cabo de 3 horas. ¿A qué velocidad va cada uno?

\ex{3} Si sumamos 1 al numerador y al denominador de una fracción, obtenemos $\tfrac23$; si restamos 1 a los dos, obtenemos $\tfrac12$. ¿Cuál es la fracción?

\ex{3} Un hotel tiene 50 habitaciones, unas dobles y otras individuales, y en total 87 camas. ¿Cuántas habitaciones hay de cada tipo?

\ex{3} El perímetro de un triángulo isósceles mide 40 cm y cada uno de los lados iguales mide 5 cm más que la base. ¿Cuánto miden los lados?

\vfill
\begin{center}\color{muted}\small\faCheckCircle[regular]\ \ ¡Ánimo! Si te cuesta algún ejercicio, repasa los apuntes interactivos de la unidad.\end{center}

\end{document}
